\section{导读：资本视角下的 AI 与互联网再定价}

本节先说明本期为什么要按“资本如何给 AI 算账”来读，而不是按投资标的逐个摘录。这样组织之后，读者可以把 OpenAI、Robinhood、Agentic Commerce 和 AI bubble 放进同一张收入/成本/入口权地图里。

本期的嘉宾 Freda 是 Altimeter Capital 合伙人，视角横跨一级和二级市场。它和 EP127 的广密季报互相补位：EP127 更像技术与产业范式复盘，本期则像资本如何给 AI、互联网金融和未来三年美股主线“算账”。节目里的数字很多，但它的教学价值不在于让读者背某个估值，而在于理解资本市场如何把技术叙事翻译成收入、现金流、市场份额、组织能力、创始人性格和二级市场表达。

本讲义不构成任何投资建议。正文保留节目中的大数和框架，是为了帮助读者理解：为什么 OpenAI 被看成产品公司而不只是模型公司；为什么模型公司现金流像“负向滚雪球”；为什么 Robinhood 一手烂牌仍可能走出 Alpha；为什么 Agentic Commerce 可能重写广告、电商和支付；以及为什么判断 AI bubble 不能只问情绪，要看模型是否继续进步、收入是否兑现、ROIC 是否改善。

\begin{importantbox}{本期核心命题}
AI 的资本叙事正在从“谁投入最多”转向“谁把投入变成收入”。资本市场不再只奖励 CAPEX 和宏大愿景，而是在追问：模型公司能否把 ChatGPT/API/Agent/广告做成收入栈，AI 应用能否从电子收入池转向劳动力池，传统行业能否把 AI 提效写进财务报表。
\end{importantbox}

\begin{knowledgebox}{视觉策略说明}
本视频是固定访谈画面，没有 slides、白板或产品演示。正文使用封面作为来源识别，正文图像全部采用简洁概念图，用来展示投资框架、收入栈、业务飞轮和泡沫判断仪表盘；细节解释放在正文和读图框中。
\end{knowledgebox}

\subsection{本章小结}

EP125 是一份资本市场观察笔记。它的主线不是“推荐买什么”，而是把 AI 和互联网公司的商业模式拆成可讨论的变量：收入从哪里来，成本什么时候放缓，谁有定价权，谁有组织执行力，谁能把趋势变成可量化财务结果。

